\chapter{Antiphospholipid Antibodies}

\section*{What Is This Test?}
Antiphospholipid antibody testing evaluates antibodies associated with antiphospholipid syndrome. The principal laboratory components are lupus anticoagulant testing, anticardiolipin IgG or IgM, and anti-beta-2-glycoprotein I IgG or IgM.

\section*{Alternative Names}
APS Panel; Anticardiolipin Antibodies; Anti-beta-2-Glycoprotein I; Lupus Anticoagulant; Thrombophilia Antibody Testing.

\section*{Why This Test Is Ordered}
Testing is considered after objectively confirmed venous or arterial thrombosis in an appropriate clinical context, recurrent pregnancy morbidity, selected unexplained late pregnancy complications, or characteristic autoimmune disease. It should not be used as indiscriminate screening of healthy people.

\section*{How This Test Is Performed}
Blood is collected for phospholipid-dependent clotting assays and immunoassays. A positive result generally requires repeat testing at least 12 weeks later to demonstrate persistence. Anticoagulants can interfere with lupus-anticoagulant testing.

\section*{Preparation, Risks, and Limitations}
No fasting is required. Risks are limited to venipuncture. Infection, pregnancy, medicines, and anticoagulant therapy can affect results. A positive antibody without a qualifying clinical event does not establish antiphospholipid syndrome.

\section*{Understanding Results}
Persistent moderate or high antibody levels, or persistent lupus anticoagulant activity, may support antiphospholipid syndrome when the clinical criteria are also met. The result must be interpreted by a clinician familiar with coagulation testing and the patient's treatment.

\section*{References}
International consensus criteria and European Alliance of Associations for Rheumatology guidance on antiphospholipid syndrome.

\clearpage
\section*{Understanding Results and Follow-up}
The laboratory report should identify the specimen, method, units, reference interval or decision limit, and whether the result is positive, negative, abnormal, or inconclusive. Reference intervals vary with age, sex, pregnancy, specimen type, assay, and laboratory; a result outside a range is not automatically a diagnosis, and a result within a range does not always exclude disease. Discuss unexpected results with the ordering clinician, who can decide whether repeat or confirmatory testing, treatment, or referral is appropriate. Do not change medicines or delay urgent care based on an isolated result.


\section*{Regulatory Status and Authorization Date}
Regulatory status applies to the specific assay, instrument, manufacturer, and intended use rather than to the test name alone. No single approval date can be assigned to this test category. For a commercial assay, verify the current product in the relevant regulator's database: the U.S. Food and Drug Administration (FDA) clearance, approval, or authorization; India's Central Drugs Standard Control Organisation (CDSCO) licence or permission; the European Union In Vitro Diagnostic Regulation (EU IVDR) CE certificate issued through the applicable conformity-assessment route; or the United Kingdom Medicines and Healthcare products Regulatory Agency (MHRA) registration and UKCA/CE status. Laboratory-developed tests may not have an individual FDA, CDSCO, EU IVDR, or MHRA approval date.
\clearpage
